% MACULAR paper draft — target: IJDAR (Springer, subscription route: no APC,
% no page charges). At submission swap \documentclass to Springer's sn-jnl
% (Overleaf template "Springer Nature LaTeX Template") — sectioning below is
% compatible. Until then this compiles with plain article anywhere.
\documentclass[11pt]{article}

\usepackage[utf8]{inputenc}
\usepackage[T1]{fontenc}
\usepackage{amsmath}
\usepackage{booktabs}
\usepackage{graphicx}
\usepackage{multirow}
\usepackage[hidelinks]{hyperref}
\usepackage[margin=2.5cm]{geometry}
\usepackage{CJKutf8} % example strings only; body text is English

\newcommand{\cer}{\textsc{cer}}
\newcommand{\EM}{\textsc{em}}
\newcommand{\hardmask}{\texttt{hard\_mask}}
\newcommand{\leace}{\textsc{leace}}

\title{Domain-Adapted OCR for Multilingual Medical Documents:\\
Transfer, Reproducibility, and the Limits of\\ Representation-Level Redaction}

\author{Sunjun Hwang}
\date{}

\begin{document}
\maketitle

\begin{abstract}
Medical forms concentrate three difficulties that general document OCR
benchmarks do not: multilingual content including CJK scripts, downstream
clinical tasks that inherit every recognition error, and text whose most
information-dense regions are exactly the ones that must not leak. We study
all three on a controlled synthetic medical-form corpus with held-out
personally identifiable information (PII) \emph{value families} and on two
real scanned-form corpora (FUNSD, XFUND). First, we show that LoRA
adaptation of a compact vision--language model (VLM) reliably improves
character error rate in every one of 28 trained-adapter$\times$language
outcomes (real-scan Japanese median \cer{} $0.846\!\to\!0.114$), that the
gain transfers from purely rendered training data to real scans --- including
to a language absent from the training corpus entirely --- and that
\emph{individual runs are not reproducible}: an identical seed on identical
code and hardware produced after-adaptation \cer{} of $0.099$ and $0.630$,
a failure invisible to training loss. We derive a practical remedy, an
$\approx$100-region post-training evaluation gate that detects the subtlest
observed divergence with $99.7\%$ reliability at $\sim$70 seconds of compute.
Second, using OCR-counterfactual layouts that defeat geometric shortcut
learning, we show the OCR error cascade into downstream extraction is real
but only measurable once shortcuts are controlled. Third, we evaluate four
redaction mechanisms --- none, structural masking, a differentiable
gate with adversarial training, and closed-form linear concept erasure
(\leace{}) --- under discrete attacks (fresh linear and nonlinear probes,
embedding inversion scored by exact match) on three VLM backbones and real
scans. Learned differentiable redaction is dominated on every backbone;
\leace{}'s guarantee holds exactly on the distribution where it is fitted
and does not transfer to unseen PII values (cross-family linear probe
$0.889$--$0.944$ against a $0.847$ majority baseline), and no mechanism
defeats a nonlinear attacker at the mechanism's own output. The evaluation
protocol --- discrete attacks, applied both at the mechanism output and after
contextualisation, under held-out PII value families --- is, to our
knowledge, the first of its kind for document-VLM region features, and each
of its three components is necessary: removing any one hides a failure this
paper reports.
\end{abstract}

\section{Introduction}\label{sec:intro}

Optical character recognition on medical documents is not a solved
sub-problem of document AI. Three properties separate a hospital intake form
from the receipts and slide decks that dominate OCR benchmarks. The content
is multilingual in a way that stresses recognisers unevenly --- Korean and
Japanese clinical forms mix Hangul, kanji/kana, Latin drug names, and dense
numerals. The output feeds automatic extraction, so recognition errors do
not merely lower a score but cascade into structured clinical data. And the
regions that matter most for the downstream task --- names, identifiers,
contact details --- are the regions that privacy regulation requires be
controlled, so any practical pipeline couples recognition with redaction.

This paper studies that coupled problem end to end and reports what we
measured, including three negative results and one reproducibility failure
that we argue is as important as any positive number. Our contributions:

\begin{enumerate}
\item \textbf{A controlled multilingual medical-form corpus with held-out
PII value families} (\S\ref{sec:corpora}). A synthetic generator renders
Korean, Japanese and English medical forms; PII values are drawn from
\emph{disjoint families} (disjoint given-name, surname, street,
organisation, phone-prefix and ID-pattern pools) assigned per split, so no
PII value seen in training recurs at evaluation. Counterfactual layouts
break the geometric shortcuts that otherwise let models answer without
reading (\S\ref{sec:cascade}). Real scanned arms use FUNSD and XFUND.
\item \textbf{Domain adaptation results with a reproducibility analysis}
(\S\ref{sec:adapt}). LoRA adaptation of Qwen2-VL-2B improves \cer{} in all
28 trained-adapter$\times$language outcomes across synthetic and real-scan
configurations, replicated on a second backbone with a $7\times$ stronger
baseline. Purely synthetic training transfers to real scans, including to a
language absent from synthetic training. But identical seeds do not pin
outcomes: the across-run distribution is heavy-tailed, divergence is
invisible to training loss, and we quantify the evaluation gate that
deployment therefore requires.
\item \textbf{A redaction mechanism comparison under a three-component
attack protocol} (\S\ref{sec:redact}): discrete attacks (probes and
embedding inversion with a shuffled-control prior floor), measured both at
the mechanism output and after relational contextualisation, under held-out
PII families. Differentiable gated redaction is dominated everywhere and
destabilises training; closed-form linear erasure is exactly right
in-distribution and vacuous under value shift; structural masking is never
worse on privacy, and its utility cost is backbone-dependent in a way we
could not predict from any cross-backbone quantity.
\end{enumerate}

Throughout we follow one reporting rule, adopted after early results failed
to replicate: \emph{a single-seed number is not a result}. Where our own
earlier conclusions were overturned by seeding, we say so explicitly
(\S\ref{sec:repro}, \S\ref{sec:discussion}).

\section{Related work}\label{sec:related}

\paragraph{Document VLMs and OCR.} Compact vision--language models such as
Qwen2-VL~\cite{wang2024qwen2vl} and its successor
Qwen2.5-VL~\cite{bai2025qwen25vl} perform OCR through free-form generation
with dynamic-resolution vision encoders; document-specialised models such as
PaddleOCR-VL~\cite{paddleocrvl} target the same setting with smaller
parameter budgets. Synthetic document generation for OCR pre-training is
established (e.g.\ SynthDoG~\cite{kim2022donut}); our generator differs in
producing \emph{semantically controlled} medical forms with per-region PII
type labels and family-disjoint value pools, which is what enables the
held-out-value evaluations in \S\ref{sec:adapt} and \S\ref{sec:redact}.

\paragraph{Parameter-efficient adaptation.} LoRA~\cite{hu2022lora} is the
standard parameter-efficient fine-tuning method. Reports of instability in
low-data fine-tuning are common lore, but we are not aware of prior work
demonstrating \emph{seed-identical} divergence (same code, data, hardware
and seed; $6\times$ \cer{} spread) for VLM OCR adaptation, nor of a sizing
analysis for the post-training acceptance gate this implies.

\paragraph{Attribute removal and its failure modes.} Adversarial attribute
removal via gradient reversal hides an attribute from the training-time
adversary without removing it; post-hoc probes recover
it~\cite{elazar2018adversarial,gonen2019lipstick} (see
\cite{barrett2019adversarial} for a counter-nuance). The successor line with
a guarantee is closed-form linear concept erasure:
INLP~\cite{ravfogel2020null}, RLACE~\cite{ravfogel2022linear}, and
\leace{}~\cite{belrose2023leace}, which provably defeats \emph{all linear}
predictors of the concept on the distribution where it is fitted. The
guarantee is linear-only, and published stress tests recover erased concepts
nonlinearly; our results quantify both this and a failure mode we have not
seen reported: collapse of the guarantee under \emph{value-family shift}
within the same concept.

\paragraph{Embedding inversion.} Cosine similarity to the original
representation is not an accepted leakage measure. The standard is
inversion: Vec2Text~\cite{morris2023vec2text} reconstructs 32-token inputs
at 92\% exact match from text embeddings, GEIA~\cite{gu2023geia} attacks
sentence embeddings generatively, and CapRecover~\cite{caprecover2025}
recovers content from intermediate \emph{vision} features --- the closest
threat model to ours. We adopt discrete scoring (exact match, \cer{}) with a
shuffled-representation prior floor, and report membership-style claims
nowhere (cf.\ LiRA~\cite{carlini2022membership} on why average-case metrics
mislead).

\paragraph{Datasets.} FUNSD~\cite{jaume2019funsd} provides real scanned
forms with question/answer region annotations; XFUND~\cite{xu2022xfund}
extends the scheme to seven languages including Japanese and Chinese. No
public Korean scanned medical-form corpus exists --- four independent
literature searches (\S\ref{sec:discussion}) found none --- which is why our
Korean results are synthetic-only and why the synthetic$\to$real transfer
result matters.

\section{Corpora and experimental protocol}\label{sec:corpora}

\paragraph{Synthetic medical forms.} A generator renders single-page medical
forms (patient headers, medication tables, vitals, free-text notes) in
Korean, Japanese and English. Every text region carries a gold transcript, a
language tag, a clinical type (from a fixed vocabulary) and a PII type
(patient name, phone, address, national ID, \dots, or non-PII).

\paragraph{Held-out PII value families.} PII values are drawn from named
families (A--E). Families have \emph{disjoint} given-name, surname, street,
organisation and phone-prefix pools, and disjoint ID century digits. Splits
use disjoint families (train A, validation B, test C; D and E exist for the
multi-family fitting experiment of \S\ref{sec:leacetransfer}), so an
evaluation-time PII value never occurred in training. Any improvement or any
leakage measured at evaluation therefore cannot be memorisation of the
literal values.

\paragraph{Counterfactual layouts.} On naturally laid-out forms a model can
learn \emph{where} answers live rather than reading them. We verify this
with a coordinate-only audit classifier (geometry, no text): F1 $\approx
1.0$ on default layouts versus $0.432$ with precision at the base rate on
counterfactual layouts, in which region positions are shuffled against
content. All reading-sensitive experiments use counterfactual layouts.

\paragraph{Real scans.} FUNSD (149 train / 50 test documents, 9{,}529
regions) supplies real scanned English forms; its \texttt{answer} regions
(filled-in values) map to the sensitive class and \texttt{question}/%
\texttt{header} regions to the utility class for the redaction experiments.
XFUND supplies real scanned Japanese, Spanish and Chinese forms for the
adaptation experiments. XFUND ships a single public split; where we must
split it ourselves we use language-stratified halves and say so --- such
splits support ``adaptation helps on these scans,'' not any unseen-value
claim.

\paragraph{Seeding policy.} Every reported experiment uses $\geq 3$ seeds;
comparisons are made only when seed-to-seed variation is reported alongside.
Where the across-run distribution turned out to be heavy-tailed
(\S\ref{sec:repro}), we pool runs and report the distribution rather than a
mean$\pm$std that would imply draws are labelled.

\section{Domain adaptation}\label{sec:adapt}

\subsection{Setup}

We LoRA-fine-tune Qwen2-VL-2B-Instruct~\cite{wang2024qwen2vl} (29M adapter
parameters; $r{=}16$, $\alpha{=}32$ unless stated; 2 epochs unless stated;
AdamW, lr $10^{-4}$, gradient clipping at 1.0) on single-region crops with
the fixed prompt ``\emph{Transcribe the text in this image exactly. Output
only the text.}'', training loss on the target tokens only. Evaluation is a
paired before/after comparison on identical crops with greedy decoding.
Crops are capped at 512\,px on the longest side; the cap is load-bearing
(uncapped real-scan crops explode into thousands of visual tokens under
dynamic resolution) and applied identically everywhere. We report
per-language \cer{} (length-weighted) and exact match (\EM{}); word error
rate is meaningless for ko/ja/zh and reported nowhere.

\subsection{Synthetic forms: everything improves}\label{sec:adapt-syn}

Training uses family A, evaluation family C (1{,}440 / 960 regions): no
shared PII values. Table~\ref{tab:syn} shows all three languages improve on
both metrics in every seed. An earlier single-seed read of this experiment
concluded English regresses; three seeds show the opposite --- one of several
single-run conclusions this project had to retract (\S\ref{sec:discussion}).

\begin{table}[t]
\centering\small
\caption{Synthetic forms, family-held-out (train A, eval C), 3 seeds.
Negative $\Delta$\cer{} = improvement.}
\label{tab:syn}
\begin{tabular}{lcccc}
\toprule
Lang & \cer{} before & \cer{} after (per seed) & $\Delta$\cer{} & $\Delta$\EM{} \\
\midrule
ko & 0.121 & 0.016 / 0.017 / 0.071 & $-0.086\pm0.026$ & $+0.084\pm0.009$ \\
ja & 0.104 & 0.015 / 0.042 / 0.073 & $-0.061\pm0.024$ & $+0.090\pm0.028$ \\
en & 0.047 & 0.006 / 0.013 / 0.036 & $-0.029\pm0.013$ & $+0.065\pm0.009$ \\
\midrule
macro & 0.088 & --- & $-0.057\pm0.020$ & --- \\
\bottomrule
\end{tabular}
\end{table}

\subsection{Real scans: large gains, language-opposite exact-match
behaviour}\label{sec:adapt-real}

On XFUND we train two configurations (ja+es and ja+zh) on
language-stratified halves. The Japanese evaluation half is \emph{identical}
in both, making the co-training language the only difference. Pooling all
trained adapters under the final code (7 per configuration;
Table~\ref{tab:real}):

\begin{table}[t]
\centering\small
\caption{Real scans (XFUND), all trained adapters pooled. Baselines: ja
0.846/0.400, es 0.924/0.520, zh 0.481/0.660 (\cer{}/\EM{}).}
\label{tab:real}
\begin{tabular}{llll}
\toprule
Config & Lang & after \cer{} (sorted) & after \EM{} (same order) \\
\midrule
ja+es & ja & 0.099 0.099 0.104 0.114 0.139 0.147 \textbf{0.630} &
             0.824 0.798 0.789 0.759 0.751 0.729 \textbf{0.100} \\
ja+es & es & 0.132 0.170 0.227 0.299 0.329 0.375 0.399 &
             0.538 0.358 0.293 0.260 0.237 0.233 0.190 \\
ja+zh & ja & 0.129 0.155 0.171 0.198 0.200 0.243 0.332 &
             0.666 0.575 0.666 0.580 0.664 0.612 0.452 \\
ja+zh & zh & 0.080 0.109 0.113 0.125 0.185 0.215 0.254 &
             0.857 0.835 0.807 0.790 0.790 0.803 0.643 \\
\bottomrule
\end{tabular}
\end{table}

Three findings survive pooling. \cer{} improves in every run, usually by a
large factor (ja median $0.846\to0.114$); even the single diverged run
(0.630, \S\ref{sec:repro}) beats its baseline. Japanese \EM{} improves
massively in 6 of 7 ja+es runs ($0.400\to0.73$--$0.82$); Spanish \EM{}
\emph{degrades} in 6 of 7 --- the direction of the exact-match effect is
language-specific and real. And the co-training partner matters: on the
identical ja evaluation half, es co-training beats zh co-training (exact
two-sided Mann--Whitney over 7 vs.\ 7 runs: \EM{} $p=0.023$, median 0.759
vs.\ 0.612; \cer{} $p=0.051$, median 0.114 vs.\ 0.198, the test narrowly
missing threshold because the es side contains the diverged run, which we do
not exclude post hoc). The direction --- Japanese helped more by Spanish
than by Chinese co-training --- is counterintuitive; we claim no mechanism,
noting only that same-script Chinese plausibly competes for shared CJK
reading capacity while Spanish does not.

\subsection{Identical seed, identical code, different
result}\label{sec:repro}

The reproducibility control repeats the ja+es configuration with seeds
$\{0,3,4,5\}$, where seed 0 \emph{exactly} repeats a previous invocation:
same code, data, machine, seed. Seed 0 produced ja after-\cer{} $0.099$ in
the first invocation and $0.630$ in the repeat, while its sibling seeds in
the same repeat landed at $0.114$--$0.147$. Training loss does not flag the
divergent run: its loss history ($0.218\to0.114$) lies inside the range of
the healthy runs ($0.208$--$0.218\to0.099$--$0.118$). The mechanism
consistent with all observations is accumulation of bf16 GPU
non-determinism over $\sim$2{,}400 sequential optimisation steps into a
heavy-tailed outcome distribution: most runs land in a good basin, an
occasional run diverges, and nothing observable during training
distinguishes them.

The consequences are practical. A seed is not a reproducibility mechanism
for this training regime; per-run numbers are not reportable facts; the
reporting unit must be the across-run distribution. And deployment requires
a post-training \emph{evaluation gate}, since training loss cannot reject a
diverged adapter.

\subsection{Sizing the evaluation gate}\label{sec:gate}

Using per-region (prediction, gold) pairs saved from every run, we bootstrap
gate sets of $n$ regions (paired sampling: one gate set, both models run on
it) and measure how often the gate correctly ranks a diverged adapter as
worse. The subtlest observed divergence (Qwen2.5-VL seed with ja \cer{}
$0.226$ vs.\ its baseline $0.118$) is detected in $96.9\%$ of 50-region
gates, $99.7\%$ of 100-region gates and $100\%$ at 200; gross divergence
($\Delta$\cer{} $0.39$) is detected with 10 regions
(Table~\ref{tab:gate}). At $0.7$\,s/region, a 100-region gate costs
$\sim$70 seconds --- negligible against a $\sim$50-minute training run, and
we argue it should be considered part of the training procedure itself.

\begin{table}[t]
\centering\small
\caption{Bootstrap detection rate of a diverged adapter vs.\ gate size $n$
(2{,}000 resamples).}
\label{tab:gate}
\begin{tabular}{lccccc}
\toprule
Detection task & $n{=}10$ & 25 & 50 & 100 & 200 \\
\midrule
Subtle ($\Delta$\cer{} 0.11, vs.\ baseline) & 0.851 & 0.919 & 0.969 & 0.997 & 1.000 \\
Gross ($\Delta$\cer{} 0.39, vs.\ sibling)   & 1.000 & 1.000 & 1.000 & 1.000 & 1.000 \\
\bottomrule
\end{tabular}
\end{table}

\subsection{Less is more: epochs and rank}\label{sec:ablate}

Two loose ends motivated ablations: training loss \emph{rose} in epoch 2 of
every synthetic seed, and $r{=}16$ was unjustified. Holding $\alpha/r=2$
fixed so that rank varies only capacity, the pattern is monotone
(Table~\ref{tab:ablate}): shrinking the budget (1 epoch) or capacity
($r{=}8$) keeps the \cer{} gain, softens the Spanish exact-match
degradation, and tightens the spread; growing capacity ($r{=}32$) makes
everything worse and wilder, its worst run landing at es \EM{} $0.098$
against a $0.520$ baseline. We therefore recommend $r{=}8$ and 1--2 epochs
as the default, with the larger settings as the ablation. (Three seeds per
variant order the spreads consistently across four variants but cannot
estimate per-variant tail rates.)

\begin{table}[t]
\centering\small
\caption{Epoch and rank ablations, XFUND ja+es, 3 seeds each. Default row
pools the 7 runs of \S\ref{sec:adapt-real}.}
\label{tab:ablate}
\begin{tabular}{lllll}
\toprule
Variant & ja \cer{} & ja \EM{} & es \cer{} & es \EM{} \\
\midrule
1 epoch, $r{=}16$ & 0.067--0.172 & 0.537--0.804 & 0.136--0.161 & 0.407--0.638 \\
$r{=}8$, 2 ep     & 0.061--0.139 & 0.651--0.806 & 0.077--0.154 & 0.468--0.667 \\
$r{=}16$, 2 ep (default) & 0.099--0.147 (+0.630) & 0.729--0.824 (+0.100) & 0.132--0.399 & 0.190--0.538 \\
$r{=}32$, 2 ep    & 0.086--0.335 & 0.271--0.784 & 0.165--0.558 & 0.098--0.577 \\
\bottomrule
\end{tabular}
\end{table}

\subsection{A second backbone}\label{sec:backbone2}

Qwen2.5-VL-3B~\cite{bai2025qwen25vl} has a $\sim$7$\times$ stronger baseline
on the same evaluation (ja \cer{} $0.118$ vs.\ $0.846$), leaving no
junk-output cliff for adaptation to repair. Both phenomena replicate: two of
three seeds roughly halve \cer{} and raise \EM{} on both languages (ja
$0.118\to0.047/0.075$, \EM{} $0.590\to0.826/0.831$), so the gain is not
weak-model repair; and one seed lands \emph{above its own baseline} (ja
\cer{} $0.226$) --- the first observed case of adaptation actively hurting
--- while its siblings are healthy. The heavy tail is not a property of one
model.

\subsection{Synthetic-to-real transfer}\label{sec:transfer}

The question that decides whether synthetic-only Korean adaptation is
meaningful: does adaptation learned on rendered forms transfer to real
scans? Training on synthetic families (no shared documents, rendering
process, or values with the evaluation corpus) and evaluating on XFUND:
Japanese \cer{} improves in 3 of 3 seeds ($1.073\to0.236/0.295/0.310$; a
ja-only pass on a larger evaluation set agrees, $0.937\to0.202$--$0.266$),
recovering most of the in-domain \cer{} gain and part of the \EM{} gain.
Spanish --- \emph{absent from the synthetic training corpus entirely} ---
improves to in-domain levels ($1.218\to0.167$--$0.254$): the adaptation is
largely a transferable reading/format skill rather than language-specific
memorisation. The reverse direction is unstable: real-scan training helps
synthetic evaluation in only 1 of 3 seeds and hurts in 2. The useful
direction, synthetic$\to$real, is the one that works; we take this as the
strongest available support for synthetic-only adaptation of languages, like
Korean, for which no public scanned corpus exists.

\section{The OCR error cascade requires shortcut
control}\label{sec:cascade}

With frozen backbone features and controlled corruption of the OCR channel,
downstream extraction on \emph{default} layouts is flat in OCR quality
(layout F1 $0.941$ at \cer{} 0 and at \cer{} 0.5): the model answers from
geometry and never reads. On counterfactual layouts the cascade appears
(F1 $0.648\to0.530$ over the same corruption). Any claim about OCR
sensitivity of a downstream task is unmeasurable without shortcut control;
we suggest the coordinate-only audit of \S\ref{sec:corpora} as a cheap
prerequisite check for any such study.

\section{Redaction under attack}\label{sec:redact}

\subsection{Mechanisms and protocol}

Four mechanisms operate on identical frozen per-region VLM features $z$:
\texttt{none}; \hardmask{} (structural: sensitive regions replaced by a
learned constant); a differentiable \texttt{gate} trained with a
gradient-reversal adversary~\cite{elazar2018adversarial}; and
\leace{}~\cite{belrose2023leace}, the closed-form affine eraser with a
provable guarantee against all linear predictors. Downstream, a relation
graph contextualises regions; we attack both $z_{\text{safe}}$ (mechanism
output) and $z_{\text{ctx}}$ (post-graph --- the representation actually
handed to consumers). \leace{}'s linear guarantee provably survives a linear
projector but not the graph's nonlinearity, which is exactly why both
surfaces must be measured.

Attacks are discrete: a fresh linear probe, a fresh 2-layer MLP probe, and
an embedding-inversion decoder scored by exact match and \cer{}, with a
shuffled-representation \emph{prior floor} subtracted (a structured-string
decoder recovers part of any identifier distribution from the text prior
alone; the floor measured $0.000$ here, but the control is what licenses the
number). Decision rules were fixed in advance: \hardmask{} is the floor any
sophisticated mechanism must beat; a linear-probe drop without a
nonlinear-probe drop is not a privacy result; and the \leace{} residual
covariance is checked on validation, not only where fitted.

\subsection{Results across three backbones}

\begin{table}[t]
\centering\small
\caption{Post-contextualisation attacks (the deployed surface), 3 seeds,
majority baseline 0.847. Seed std 0.002--0.023, smaller than every
difference shown.}
\label{tab:redact}
\begin{tabular}{llcccc}
\toprule
Backbone & Mechanism & clinical F1 & linear & nonlinear & inv.\ \cer{} \\
\midrule
PaddleOCR-VL 0.9B & none      & 0.976 & 0.958 & 0.975 & 0.454 \\
                  & \hardmask & 0.947 & 0.849 & \textbf{0.843} & \textbf{1.080} \\
                  & gate      & 0.917 & 0.919 & 0.937 & 0.705 \\
                  & \leace    & 0.942 & 0.871 & 0.917 & 0.549 \\
\midrule
Qwen2-VL-2B       & none      & 0.899 & 0.956 & 0.968 & 0.549 \\
                  & \hardmask & 0.778 & 0.848 & \textbf{0.859} & \textbf{0.949} \\
                  & gate      & 0.826 & 0.923 & 0.929 & 0.743 \\
                  & \leace    & \textbf{0.890} & 0.846 & 0.860 & 0.844 \\
\midrule
Ministral-3       & none      & 0.944 & 0.953 & 0.975 & 0.493 \\
                  & \hardmask & 0.731 & 0.876 & \textbf{0.879} & \textbf{0.915} \\
                  & gate      & 0.779 & 0.946 & 0.959 & 0.581 \\
                  & \leace    & \textbf{0.897} & 0.877 & 0.896 & 0.623 \\
\bottomrule
\end{tabular}
\end{table}

Four claims hold on all three backbones (Table~\ref{tab:redact}), with seed
variation below every difference quoted.

\textbf{(1) Learned differentiable redaction is dominated everywhere.} On
every backbone some simpler mechanism beats the gate on utility and leakage
simultaneously; on Ministral-3 it is the worst-protected mechanism measured
while also costing utility. The gate additionally destabilises training:
across a LoRA ablation, every variant containing the gate had clinical-F1
seed std 0.37--0.41 versus 0.007 without it, because the gate makes the
downstream representation a function of a PII head that is random early in
training.

\textbf{(2) \leace{} does not transfer to unseen PII values}
(\S\ref{sec:leacetransfer}).

\textbf{(3) No mechanism defeats a nonlinear attacker at its own output}
(pre-graph nonlinear probe $\geq 0.942$ everywhere against the 0.847
baseline). Contextual mixing then pushes \hardmask{} to complete protection
(0.843, at the baseline within noise) while leaving the gate's and
\leace{}'s soft residue readable. A privacy number reported on the
mechanism's own output describes a representation nobody deploys ---
pre-graph, \hardmask{} looks partial (0.966) when it is in fact complete.

\textbf{(4) Structural masking is never worse on privacy; its utility cost
is what varies} --- and varies too much to predict. Against \leace{}, the
clinical-F1 gap ranges from $-0.005$ (PaddleOCR-VL) to $+0.166$
(Ministral-3). Value reconstruction is unambiguous: inversion exact match
(prior floor subtracted) is lowest for \hardmask{} on all four datasets
including real scans (0.013/0.002/0.020/0.008), while \leace{} leaves the
literal value nearly as recoverable as no protection on three of four
(e.g.\ 0.144 vs.\ 0.157 on PaddleOCR-VL; 0.101 vs.\ 0.109 on FUNSD). We
attempted a cross-backbone predictor of the masking-vs-erasure trade and
report its failure in \S\ref{sec:discussion}.

\subsection{\leace{} under value shift}\label{sec:leacetransfer}

Fitted on family A and applied to family B --- before any projector or
training --- the linear probe on erased features stays far above the 0.847
majority baseline on all three backbones: 0.932 / 0.889 / 0.944, against raw
0.964 / 0.975 / 0.967. The implementation is not at fault: fitting the
eraser on validation itself drives the probe to exactly 0.847 with residual
covariance $1.3\times10^{-7}$. The guarantee holds precisely where it is
fitted and nowhere else, and the deployment condition --- new patients,
unseen names --- is ``nowhere else'' by construction.

None of the cheap remedies work: coarsening the concept to binary is-PII
(0.956), covariance shrinkage (0.928), truncating the concept subspace to
its top directions (0.952--0.956) are no better than the 0.932 default.
Neither does data: fitting across three value families (17{,}176 regions,
$3\times$ the default) moves the probe from 0.932 to 0.925, closing 0.007
of the 0.085 gap. The trend is real and far too shallow to matter, making
this an intrinsic limitation rather than a data-budget problem: the linear
subspace carrying one family's identifiers is not the subspace carrying the
next family's.

Even in the best case --- eraser fitted on the evaluation distribution
itself --- a nonlinear probe retains 0.912 against the 0.847 baseline: the
linear-only caveat of the erasure literature, measured here on document-VLM
region features.

\section{Discussion}\label{sec:discussion}

\paragraph{What we recommend to practitioners.} For recognition: LoRA-adapt
a compact VLM with $r{=}8$ for 1--2 epochs, and treat a $\sim$100-region
evaluation gate as part of training --- accept the adapter only if it beats
the base model on the gate set. For redaction: structurally mask sensitive
regions before contextualisation; do not rely on learned gating (dominated,
destabilising) or linear erasure (vacuous under value shift) for anything
that must withstand reconstruction of literal values.

\paragraph{Negative results and retractions.} This project retracted nine of
its own confidently-stated intermediate conclusions, all for one of two
reasons: single-seed (or single-run) reads of heavy-tailed distributions, or
metrics that do not measure what they claim (cosine leakage; a
Frobenius-norm ``transfer quality'' that is incomparable across backbones
and failed as a predictor on the third backbone). We enumerate them in the
project's findings log rather than here, but two deserve print because they
bracket the same lesson: from single runs we first concluded real-scan
adaptation ``teaches reading'' (CER and EM moved together), then ---
from two other single runs --- that it is ``format control, not reading''
(EM fell while CER dropped). Pooled over 14 runs, neither slogan is right:
Japanese EM improves massively in 6 of 7 runs, Spanish EM degrades in 6 of
7, and both single-run reads were tail draws. The distribution is the
result; a run is an anecdote.

\paragraph{Limitations.} Korean and the family-held-out guarantees are
synthetic-only; no public Korean scanned medical corpus exists (the closest
resources are Chinese: CHIP2022 MedOCR, CC-BY-NC-SA), and we identify this
as a field gap rather than pretend around it. XFUND self-splits support
in-corpus claims only. The redaction study uses frozen features, three
backbones and 120-document splits; the adaptation study uses two backbones
of one family (Qwen2/2.5-VL). Tail rates per configuration are bounded by
our run counts (one diverged run in 14 at $r{=}16$); estimating them
precisely would take an order of magnitude more runs.

\section{Conclusion}

On a controlled multilingual medical-form corpus with held-out PII value
families, plus real scanned forms, we showed: domain adaptation of compact
VLM OCR works and transfers from synthetic to real scans (including
cross-language), but its per-run outcomes are heavy-tailed to the point that
identical seeds do not reproduce --- so distribution-level reporting and a
cheap post-training evaluation gate, which we size at $\sim$100 regions, are
not optional; the OCR error cascade into downstream extraction is real but
only measurable under shortcut control; and of four redaction mechanisms
attacked with discrete probes and inversion on two surfaces, structural
masking is the only one whose protection survives every attack, learned
gating is dominated everywhere, and closed-form linear erasure --- exactly
right where fitted --- is vacuous under the value shift that defines
deployment. The three-component evaluation protocol (discrete attacks, both
surfaces, held-out value families) is reusable as-is, and each component is
necessary: dropping any one of them would have hidden a failure this paper
reports.

\bibliographystyle{plain}
\bibliography{references}

\end{document}
